% !TEX root = ../main.tex
\chapter{Conclusion}
\label{chapter:Conclusion}

\section{Summary}
\label{section:ConclusionSummary}

This thesis set out to advance the differentiable modular planning stack developed on the RoboRacer platform.
The goals were twofold: to enable on-board execution of DiffPlan, and to broaden the family of differentiable local motion planners available within the MIND ecosystem \cite{jahncke_MINDStackModularInterpretable_2025, swadiryus_DifferentiableMotionPlanning_2025, koch_DifferentiableInterpretablePath_2025}.
Chapter~\ref{chapter:optimizations} identified the runtime bottlenecks of the original DiffPlan implementation through profiling.
Two routines were then vectorized to remove them: the longitudinal sampling routine, and the wall-proximity computation used inside the collision loss.
GPU offloading was also explored to leverage the accelerator present on the target Jetson Orin Nano Super compute platform.
Chapter~\ref{chapter:SpatialSampler} introduced DiffSpatialPlan, a differentiable spatial sampling planner that adapts the online velocity-profile method of Langmann et al.\ \cite{langmann_OnlineVelocityProfile_2025} to the two-dimensional RoboRacer setting.
DiffSpatialPlan pairs online apex detection with a forward-backward friction-limited solver and a quintic spatial parameterization.
Hard $\arg\min$ operations are replaced with soft aggregations so that gradients from downstream control losses can still propagate through the planner.

Subsequently, Chapters~\ref{chapter:Experiments} and~\ref{chapter:Results} evaluated the two contributions against the baseline planner.
The benchmarks cover runtime, scalability, and driving behavior in the F1TENTH Gym simulator, and extend to the physical RoboRacer vehicle on the FTM Oval circuit.
The vectorized DiffPlan achieved an approximately $2.6\times$ end-to-end speedup relative to the original implementation, raising the $\texttt{/drive}$ publish rate from $13.2$\,Hz to $34.1$\,Hz.
DiffSpatialPlan reached $30.7$\,Hz on the same workstation.
The scaling study confirmed that the vectorized variants soften the dependence of planning time on the number of sampled trajectories that characterized the baseline.
In the head-to-head driving comparison, DiffPlan consistently produced tighter raceline adherence, while DiffSpatialPlan produced faster and more robust laps.
The robustness gap was most notable on the boundary-constrained Austria circuit, where DiffSpatialPlan avoided the wall contacts observed with the temporal planner during early training laps.
The physical FTM Oval runs showed that both planners train successfully on hardware under motion-capture localization.
The sim-to-real gap manifested as a hardware-imposed speed ceiling that compressed the lap-time advantage previously observed for the spatial planner in simulation.

\section{Outlook}
\label{section:ConclusionOutlook}

Taken together, these results support two broader observations about the differentiable modular stack.
The first observation concerns where the runtime cost of differentiability actually comes from.
Much of that cost does not originate from automatic differentiation itself, but from the unvectorized Python control flow that surrounds the tensor operations.
Once the sampling and collision-loss routines were reformulated as batched linear-algebra expressions, the planner reached publish rates that approach the real-time regime targeted by the MIND stack \cite{jahncke_MINDStackModularInterpretable_2025}.

The second observation concerns hardware placement.
Partial migration of the stack to the GPU can be actively harmful.
Because the differentiable MPC still solves its LQR sub-problem sequentially on the CPU, tensors must be transferred from GPU to CPU and back again at every planning cycle.
In the GPU-offloading experiment, these host-device transfers erased the gains obtained elsewhere on the graph.
Differentiable modular stacks must therefore be considered as a whole when reasoning about hardware placement, since optimizing each module in isolation is not sufficient.

Moreover, the comparison between DiffPlan and DiffSpatialPlan clarifies the design trade-off between temporal and spatial sampling in a racing context.
The temporal sampler discretizes the horizon in time, so its lookahead footprint contracts at low speeds and expands at high speeds.
This tightly couples raceline adherence to the current velocity of the vehicle.
The spatial sampler instead fixes the horizon in arc length and derives velocity from an online friction-limited profile that is aware of upcoming apexes.
As a result, DiffSpatialPlan exhibits behavior that is closer to that of a human racing driver.
Braking is placed relative to the geometry of the corner rather than relative to the current sampling horizon, and cornering speed reflects the currently available grip.
This structural difference explains both the robustness advantage on Austria and the higher cumulative lateral deviation on all three tracks.
The wider cornering lines that the spatial sampler favors are geometrically further from the offline reference, even when they are dynamically preferable.

Nevertheless, one central objective of the thesis remained unmet.
Full on-board execution of the stack on the Jetson Orin Nano Super was not achieved.
The main reason is that the CPU-bound DiffMPC controller prevents the planning graph from remaining resident on the accelerator.
The algorithm therefore continues to run on the external workstation and communicates with the vehicle over ROS2, as illustrated in Figure~\ref{fig:exp_setup_real}.
In light of this limitation, the contribution of this thesis should be understood as a decisive step toward on-vehicle deployment rather than a completion of it.
The algorithmic footprint has been reduced to a level at which on-board computation becomes plausible, and the sim-to-real transferability of the training procedure has been validated on the physical FTM Oval circuit.

\section{Future Work}
\label{section:ConclusionFutureWork}

Building on these findings, several directions appear particularly promising for subsequent work on the differentiable RoboRacer stack.
Foremost, replacing the current DiffMPC controller with a fully parallelizable formulation is the most immediate prerequisite for on-vehicle deployment.
The recent GPU-accelerated differentiable MPC of Adabag et al.\ \cite{adabag_DifferentiableModelPredictive_2025}, implemented in JAX, offers a natural candidate.
Migrating the controller in this direction would eliminate the host-device oscillation identified in Section~\ref{section:GPUOffloading} and allow the entire planning-control graph to remain resident on the Jetson GPU.
Closing the controller gap alone should suffice to reach onboard real-time execution.

In addition, extending differentiability toward a more complex prediction module constitutes a natural continuation of the MIND research programme \cite{jahncke_MINDStackModularInterpretable_2025, huang_DifferentiableIntegratedMotion_2024, karkus_DiffStackDifferentiableModular_2022}.
The experiments in this thesis intentionally isolated the planning and control stack.
They relied on ground-truth pose from simulator odometry or motion capture, and the prediction module was disabled.
Reintroducing a differentiable opponent-prediction model, together with the differentiable perception front-end, would allow downstream driving losses to shape upstream feature extraction.
Care must be taken to select architectures whose gradients remain informative under the soft-selection mechanisms already used by DiffPlan and DiffSpatialPlan.

Furthermore, the evaluation methodology itself can be strengthened.
The driving experiments in this thesis were limited to solo lap-time optimization without opponents or static obstacles, and to a single physical circuit.
Follow-up work should evaluate DiffSpatialPlan on additional physical tracks, incorporate opponent vehicles to exercise the currently dormant prediction path, and extend the sim-to-real study observed on the FTM Oval circuit.
Ultimately, three ingredients would move the RoboRacer platform closer to the vision of an interpretable, gradient-trainable autonomous racing stack that operates entirely on-board: a parallelizable differentiable controller, a fully differentiable perception-prediction-planning-control pipeline, and broader empirical validation.
